% ACTL3142 Week 1 Tutorial Slides
\documentclass[aspectratio=169]{beamer}

% Add beamer theme directory to search paths
\makeatletter
\providecommand*{\input@path}{}
\g@addto@macro\input@path{{../beamer-unsw-theme-master/}}
\makeatother
\graphicspath{{../beamer-unsw-theme-master/}{assets/}}

\usetheme[
	logo=images/unsw-portrait.png,
	sidelogo=images/unsw-landscape.png,
]{unsw}

\usepackage{amsmath,amssymb}
\usepackage{booktabs}
\usepackage{hyperref}

% information for the title page
\author{\texorpdfstring{Nick Guo\\{\small With references to Tadhg Xu-Glassop's slides}}{Nick Guo}}
\title{ACTL3142 Tutorial}
\subtitle{Week 1: Introduction to Statistical Learning}
\institute{School of Risk and Actuarial Studies, UNSW Sydney}
\date{T1 2026}

\begin{document}

% ============================================================
% TITLE
% ============================================================
\begin{frame}[plain]
	\titlepage
\end{frame}

% ============================================================
% INTRODUCTION
% ============================================================
\begin{frame}{Introduction}
	Hi! My name is \textbf{Nick} and I'm excited to be tutoring you this term.

	\bigskip
	\textbf{A little about me:}
	\begin{itemize}
		\item Fourth-year Actuarial Studies and Computer Science student
		\item Interested in poker and basketball
		\item I took this course in T2 2024
	\end{itemize}

	\bigskip
	\begin{alertblock}{Advice}
		This course is far too rich to learn exclusively within tutorials. Please keep up with lectures and ideally read the website and textbook. The exam is typed (Inspera), so focus on \textbf{intuition and application} rather than memorising formulas.
	\end{alertblock}
\end{frame}

% ============================================================
% ICEBREAKERS
% ============================================================
\begin{frame}{Icebreakers}
	Please go around and introduce yourself:

	\bigskip
	\begin{enumerate}
		\item What's your name?
		\item Year and degree
		\item Favourite course you've taken
		\item An emoji to describe your day so far
	\end{enumerate}
\end{frame}

% ============================================================
% FUNDAMENTALS OF STATISTICAL LEARNING
% ============================================================
\section{Fundamentals of Statistical Learning}

\begin{frame}{What is Statistical Learning?}
	A vast set of tools for \textbf{understanding data} and \textbf{predicting outcomes}.

	\bigskip
	We believe that data follows a model of the form:
	\[
		Y = f(\mathbf{X}) + \varepsilon
	\]
	\begin{itemize}
		\item $f$ is some unknown function that captures the \textbf{systematic} (deterministic) relationship between inputs and output
		\item $\mathbf{X} = (X_1, \dots, X_p)$: the predictors (features, independent variables)
		\item $Y$: the response (output, dependent variable)
		\item $\varepsilon$: the irreducible error, i.e.\ everything $f(\mathbf{X})$ cannot explain (random chance, measurement error, missing variables, etc.)
	\end{itemize}

	\medskip
	Our goal: estimate $f$. The resulting $\hat{f}$ is our \textbf{model}.
\end{frame}

\begin{frame}{Prediction vs.\ Inference}
	\begin{columns}[T]
		\begin{column}{0.48\textwidth}
			\textbf{Prediction}
			\begin{itemize}
				\item Goal: accurately estimate $Y$ given $\mathbf{X}$
				\item The form of $f$ is not the focus; we just want accurate outputs
				\item Models can be very complex (black box is fine)
			\end{itemize}
		\end{column}
		\begin{column}{0.48\textwidth}
			\textbf{Inference}
			\begin{itemize}
				\item Goal: understand \emph{how} $\mathbf{X}$ affects $Y$
				\item Which predictors matter? What is their individual effect?
				\item Models tend to be simpler and more interpretable
			\end{itemize}
		\end{column}
	\end{columns}

	\bigskip
	\begin{exampleblock}{Key questions to always ask}
		What predictors should we include? Do we want to predict or explain? What is signal versus noise?
	\end{exampleblock}
\end{frame}

\begin{frame}{Key Distinctions}
	\begin{itemize}
		\item \textbf{Supervised} ($Y$ observed) vs.\ \textbf{Unsupervised} (no $Y$; find structure)
		\item \textbf{Regression} ($Y$ quantitative) vs.\ \textbf{Classification} ($Y$ categorical)
		\item \textbf{Interpretability} vs.\ \textbf{Flexibility}: easy to understand $f$ \emph{or} able to capture complex relationships (typically a tradeoff)
	\end{itemize}

	\bigskip
	\textbf{Two approaches to estimating $f$:}
	\begin{columns}[T]
		\begin{column}{0.48\textwidth}
			\emph{Parametric}
			\begin{itemize}
				\item Assume a form for $f$ (e.g.\ linear)
				\item Estimate a few parameters
				\item Works with limited data
				\item Risk: misses signal (\alert{high bias})
			\end{itemize}
		\end{column}
		\begin{column}{0.48\textwidth}
			\emph{Non-parametric}
			\begin{itemize}
				\item No assumption about $f$'s shape
				\item Many effective ``parameters''
				\item Needs lots of data
				\item Risk: fits noise (\alert{high variance})
			\end{itemize}
		\end{column}
	\end{columns}
\end{frame}

% ============================================================
% CQ2
% ============================================================
\section{Conceptual Questions}

\begin{frame}{CQ2: Classification or Regression? (ISLR2 Q2.2)}
	\begin{center}
		\includegraphics[width=0.85\textwidth]{CQ2}
	\end{center}
\end{frame}

% ============================================================
% MODEL FITTING, BIAS-VARIANCE TRADEOFF, AND METRICS
% ============================================================
\section{Model Fitting and the Bias-Variance Tradeoff}

\begin{frame}{Loss Functions and the Bias-Variance Tradeoff}
	Suppose our model predicts $\hat{y}$. We compare this to the true $y$ using a \textbf{loss function} (e.g.\ squared error).

	\bigskip
	The expected test MSE decomposes as:
	\[
		\underbrace{E\!\big[(y_0 - \hat{f}(x_0))^2\big]}_{\text{Expected Test MSE}}
		\;=\;
		\underbrace{\big[\mathrm{Bias}\!\big(\hat{f}(x_0)\big)\big]^2}_{\text{Bias}^2}
		\;+\;
		\underbrace{\mathrm{Var}\!\big(\hat{f}(x_0)\big)}_{\text{Variance}}
		\;+\;
		\underbrace{\mathrm{Var}(\varepsilon)}_{\text{Irreducible}}
	\]

	\begin{itemize}
		\item \textbf{Bias}: systematic error from simplifying assumptions
		\item \textbf{Variance}: how much $\hat{f}$ changes across different training sets
		\item $\mathrm{Var}(\varepsilon)$: irreducible error, a floor we cannot beat
	\end{itemize}
\end{frame}

\begin{frame}{Visualising the Tradeoff}
	\begin{center}
		\includegraphics[width=0.55\textwidth]{bv_tradeoff}
	\end{center}

	\begin{itemize}
		\item More flexible $\Rightarrow$ less bias, more variance
		\item Less flexible $\Rightarrow$ more bias, less variance
		\item The best model achieves an \textbf{optimal tradeoff} at the minimum of total error
	\end{itemize}
\end{frame}

\begin{frame}{Bias-Variance in Action: KNN}
	\textbf{KNN classifier}: to classify a new point, find its $K$ nearest neighbours and let them vote. The majority class wins.

	\smallskip
	\begin{center}
		\includegraphics[width=0.50\textwidth]{knn}
	\end{center}

	\smallskip
	\begin{columns}[T]
		\begin{column}{0.48\textwidth}
			\textbf{$K = 1$} (very flexible)\\
			Low bias, high variance $\to$ \alert{overfitting}
		\end{column}
		\begin{column}{0.48\textwidth}
			\textbf{$K = 100$} (inflexible)\\
			High bias, low variance $\to$ \alert{underfitting}
		\end{column}
	\end{columns}

	\medskip
	\noindent Some $K$ between 1 and 100 achieves a better bias-variance tradeoff.
\end{frame}

\begin{frame}{Simple Linear Regression: Let's Work Through an Example}
	A \textbf{parametric} approach. We assume a structure for $f$:
	\[
		y = \beta_0 + \beta_1 x + \varepsilon
	\]

	\textbf{Interpreting the coefficients:}
	\begin{itemize}
		\item $\beta_0$ (intercept): the predicted value of $y$ when $x = 0$
		\item $\beta_1$ (slope): the average change in $y$ for a one-unit increase in $x$
	\end{itemize}

	\bigskip
	\textbf{Properties:}
	\begin{itemize}
		\item Not very flexible, since we impose a linearity assumption
		\item But easy to interpret: a clear example of the flexibility-interpretability tradeoff
	\end{itemize}

	\bigskip
	\textbf{Estimating the coefficients:}
	\begin{itemize}
		\item Given training data, we need assumptions about $\varepsilon$ to estimate $\beta_0$ and $\beta_1$
		\item The idea: pick the $\beta$'s that \textbf{minimise the loss function}
		\item Here we use \textbf{MSE} (Mean Squared Error) as our loss
	\end{itemize}
\end{frame}

% ============================================================
% REMAINING CONCEPTUAL QUESTIONS
% ============================================================

\begin{frame}{CQ1: Flexible vs.\ Inflexible Methods (ISLR2 Q2.1)}
	\begin{center}
		\includegraphics[width=0.85\textwidth]{CQ1}
	\end{center}
\end{frame}

\begin{frame}{CQ3: Bias-Variance Decomposition (ISLR2 Q2.3)}
	\begin{center}
		\includegraphics[width=0.85\textwidth]{CQ3}
	\end{center}
\end{frame}

\begin{frame}{CQ4: Advantages and Disadvantages of Flexibility (ISLR2 Q2.5)}
	\begin{center}
		\includegraphics[width=0.85\textwidth]{CQ4}
	\end{center}
\end{frame}

% ============================================================
% APPLIED QUESTIONS
% ============================================================
\section{Applied Questions}

\begin{frame}{Applied Questions}
	Time to switch to R and work through the applied exercises.

	\bigskip
	\textbf{What we'll cover:}
	\begin{itemize}
		\item Exploratory data analysis (EDA) with real datasets
		\item Using \texttt{summary()}, \texttt{pairs()}, \texttt{plot()}, and \texttt{hist()}
		\item Understanding your data \emph{before} fitting any models
	\end{itemize}

	\bigskip
	\begin{exampleblock}{Key Takeaway}
		Always explore your data first. Summary statistics and visualisations reveal patterns, outliers, and potential issues that inform every modelling decision you make.
	\end{exampleblock}
\end{frame}

\begin{frame}{Thanks for coming! See you next week.}
	\begin{center}
		\includegraphics[height=0.75\textheight]{coco}
	\end{center}
\end{frame}

\end{document}
